\documentclass[a4paper,12pt]{article}
\usepackage{amssymb,amsmath,amsfonts,amsthm}
\usepackage{graphicx}
\usepackage{fullpage}
\usepackage{multirow}
\usepackage{colortbl}
\usepackage[latin1]{inputenc}
\usepackage[T1]{fontenc}
\usepackage[brazil]{babel}
\usepackage{calligra}
\usepackage{accents}
\usepackage{cancel}

% Definindo o estilo de numera��o dos subteoremas
\newtheorem{theorem}{Teorema}[section]

% Definindo um novo estilo espec�fico para subteoremas com a numera��o desejada
\newtheorem{subtheorem}{Teorema}[theorem]

% Customizando o contador do subtheorem
\renewcommand{\thesubtheorem}{\thesection.\arabic{theorem}(\alph{subtheorem})}

\begin{document}
	
% Define a cor da linha de cancelamento como vermelha
\renewcommand{\CancelColor}{\color{red}}

\setcounter{section}{13}

\setcounter{theorem}{1}

\setcounter{subtheorem}{0}

\begin{subtheorem}
	Se $X\sim \text{Hipergeom�trica}(N,n,k)$, a m�dia  � dada por $ E(X) = \frac{nk}{N} $.
\end{subtheorem}

\begin{proof}
	Temos que
	$$
		\begin{aligned}
			E(X) &= \sum\limits_{x=0}^n x P(X=x)\\
			&= \sum\limits_{x=0}^n x \frac{\binom{k}{x}\binom{N-k}{n-x}}{\binom{N}{n}}\\
			&= 0 + \sum\limits_{x=1}^n x \frac{\binom{k}{x}\binom{N-k}{n-x}}{\binom{N}{n}}\\
			&= \sum\limits_{x=1}^n x \frac{ \frac{k!}{x!(k-x)!} \binom{N-k}{n-x} }{ \frac{N!}{n!(N-n)!} }\\
			&= \sum\limits_{x=1}^n x \frac{ \frac{k(k-1)!}{x(x-1)!(k-x)!} \binom{N-k}{n-x} }{ \frac{N(N-1)!}{n(n-1)!(N-n)!} }\\
			&= \sum\limits_{x=1}^n x \frac{\frac{k}{x}}{\frac{N}{n}} \frac{ \frac{(k-1)!}{(x-1)!(k-x)!} \binom{N-k}{n-x} }{ \frac{(N-1)!}{(n-1)!(N-n)!} }\\
			&= \sum\limits_{x=1}^n \cancel{x} \frac{k n}{\cancel{x} N} \frac{ \binom{k-1}{x-1} \binom{N-k}{n-x} }{ \binom{N-1}{n-1} }\\
			&= \frac{k n}{N} \sum\limits_{x=1}^n \frac{ \binom{k-1}{x-1} \binom{N-k}{n-x} }{ \binom{N-1}{n-1} }.
		\end{aligned}
	$$
	Fazendo $y=x-1$, temos
	$$
		E(X) = \frac{k n}{N} \sum\limits_{y=0}^{n-1} \frac{ \binom{k-1}{y} \binom{N-k}{n-(y+1)} }{ \binom{N-1}{n-1} }.
	$$
	Agora fazendo $M = N-1$, $m = n-1$ e $c = k-1$ no somat�rio, temos
	$$
		E(X) = \frac{k n}{N} \sum\limits_{y=0}^m \frac{ \binom{c}{y} \binom{M-c}{m-y} }{ \binom{M}{m} } = \frac{k n}{N} \sum\limits_{y=0}^m P(Y=y),
	$$
	em que $Y\sim \text{Hipergeom�trica}(M,m,c)$. Logo, $\sum\limits_{y=0}^m P(Y=y) = 1$ e
	$$
		E(X) = \frac{k n}{N}.
	$$
\end{proof}

\begin{subtheorem}
	Se $X\sim \text{Hipergeom�trica}(N,n,k)$, a vari�ncia  � dada por $$ Var(X) = n\frac{k}{N} \frac{(N-k)}{N} \frac{N-n}{N-1} $$.
\end{subtheorem}
	
\begin{proof}
	Sabemos que a vari�ncia � definida como $Var(X) = E(X^2) - [E(X)]^2$. Vamos calcular primeiro $E(X^2)$.
	$$
	\begin{aligned}
		E(X^2) &= \sum\limits_{x=0}^n x^2 P(X=x)\\
		&= \sum\limits_{x=0}^n x^2 \frac{\binom{k}{x}\binom{N-k}{n-x}}{\binom{N}{n}}\\
		&= 0 + \sum\limits_{x=1}^n x^2 \frac{\binom{k}{x}\binom{N-k}{n-x}}{\binom{N}{n}}\\
		&= \sum\limits_{x=1}^n x\left[ x \frac{\binom{k}{x}\binom{N-k}{n-x}}{\binom{N}{n}} \right]\\
		&= \sum\limits_{x=1}^n x \left[ \frac{kn}{N} \frac{ \binom{k-1}{x-1} \binom{N-k}{n-x} }{ \binom{N-1}{n-1} } \right]\\
		&= \frac{kn}{N} \sum\limits_{x=1}^n x \frac{ \binom{k-1}{x-1} \binom{N-k}{n-x} }{ \binom{N-1}{n-1} }\\
	\end{aligned}
	$$
	Fazendo $y=x-1$, temos
	$$
		E(X^2) = \frac{kn}{N} \sum\limits_{y=0}^{n-1} (y+1) \frac{ \binom{k-1}{y} \binom{N-k}{n-(y+1)} }{ \binom{N-1}{n-1} }.
	$$
	Fazendo agora $M = N-1$, $m = n-1$ e $c = k-1$ no somat�rio, temos
	$$
	\begin{aligned}
		E(X^2) &= \frac{kn}{N} \sum\limits_{y=0}^{m} (y+1) \frac{ \binom{c}{y} \binom{M-c}{m-y} }{ \binom{M}{m} }\\
		&= \frac{kn}{N} \sum\limits_{y=0}^{m} \left[ y \frac{ \binom{c}{y} \binom{M-c}{m-y} }{ \binom{M}{m} } + \frac{ \binom{c}{y} \binom{M-c}{m-y} }{ \binom{M}{m} } \right]\\
		&= \frac{kn}{N} \left[ \sum\limits_{y=0}^{m} y \frac{ \binom{c}{y} \binom{M-c}{m-y} }{ \binom{M}{m} } + \sum\limits_{y=0}^{m} \frac{ \binom{c}{y} \binom{M-c}{m-y} }{ \binom{M}{m} } \right]\\
		&= \frac{kn}{N} \left[ \sum\limits_{y=0}^{m} y P(Y=y) + \sum\limits_{y=0}^{m} P(Y=y) \right],
	\end{aligned}
	$$
	em que $Y\sim \text{Hipergeom�trica}(M,m,c)$. Logo, $$\sum\limits_{y=0}^{m} y P(Y=y) = E(Y) = \frac{cm}{M} \quad \text{e} \quad \sum\limits_{y=0}^n P(Y=y) = 1.$$
	Portanto,
	$$
		\begin{aligned}
			E(X^2) &= \frac{kn}{N} \left[ \frac{cm}{M} + 1 \right]\\
			&= \frac{kn}{N} \left[ \frac{(k-1)(n-1)}{N-1} + 1 \right]
		\end{aligned}
	$$
	Assim,
	$$
		\begin{aligned}
			Var(X) &= E(X^2) - [E(X)]^2\\
			&= \frac{kn}{N} \left[ \frac{(k-1)(n-1)}{N-1} + 1 \right] - \left[ \frac{k n}{N} \right]^2\\
			&= \frac{kn}{N} \left[ \frac{kn-k-n+1}{N-1} + 1 - \frac{k n}{N}\right]\\
			&= \frac{kn}{N} \left[ \frac{N(kn-k-n+1) + N(N-1) - kn(N-1)}{N(N-1)} \right]\\
			&= \frac{kn}{N} \left[ \frac{\cancel{Nkn} -Nk -Nn +\bcancel{N} + N^2 -\bcancel{N} - \cancel{Nkn} + kn)}{N(N-1)} \right]\\
			&= \frac{kn}{N} \left[ \frac{N^2 -Nn -Nk +kn}{N(N-1)} \right]\\
			&= \frac{kn}{N} \frac{(N-k)(N-n)}{N(N-1)}\\
			&= n\frac{k}{N} \frac{(N-k)}{N} \frac{N-n}{N-1}
		\end{aligned}
	$$
\end{proof}	
	
	
	
	
	


\end{document}